
\section{What are Zero-Knowledge Proofs?}

All of us have heard about Zero-Knowledge, but what is it exactly? In my experience, many projects have abused this term --- using it only to hide something already known, like multi-signature protocols, or centralising data and value behind a web2 backend while claiming ZK legitimacy.

A telling example: a project once claimed to make all transactions anonymous. You would send stablecoins to them, tell them the intended receiver, and they would forward the funds bundled with other transfers so the origin was obscured. There is no ZK in that process --- it is simply a centralised mixer.

\subsection*{Where's Waldo and the Hash Lock Analogy}

The example that first made ZK click for me is \textbf{Where's Waldo} (known as \textit{Where's Wally} outside North America). Imagine you have found Waldo on a busy page but you want to prove to a friend that you know where he is without pointing at him. You could take a very large sheet of paper, cut a small hole over Waldo, and show just his face through the hole. Your friend sees Waldo and is convinced you know the location --- without learning anything else about the page.

Using secrets in blockchains was already possible in Bitcoin with \textbf{hash-lock scripts}: you choose a password known only to you, hash it, and lock funds to that hash. Anyone who can reveal the pre-image unlocks the coins.

\begin{remark}
Hash Time-Locked Contracts (HTLCs) are the technology behind cross-chain atomic swaps --- used by protocols such as \textit{Fluid Tokens} --- and underpins payment routing on the Bitcoin Lightning Network.
\end{remark}

The critical limitation of hash locks is that revealing the password burns it: once the pre-image is public, that address is permanently unsafe because anyone now knows the secret tied to it. With ZK you can \textbf{prove that you know the password without ever revealing it}. That is the real magic.

\subsection*{ZK as Compression, Not Just Privacy}

Most people stop at the privacy angle. ZK is far more useful as a \textbf{compression tool}, and this is arguably its most important application in blockchain today.

\paragraph{ZK Rollups (general case).}
An L2 network may process thousands of transactions per second. Posting every transaction to L1 would be prohibitively expensive. Instead, the rollup posts only a ZK proof attesting that all those transactions were correctly signed by the users involved. An L1 smart contract verifies the proof --- a single compact check replaces thousands of individual verifications.

\begin{highlight}
\textbf{Thousands of L2 transactions $\longrightarrow$ one ZK proof on L1.}\\
The on-chain contract only verifies the proof. Nobody can lie about what happened.
\end{highlight}

\paragraph{The Cardano use case.}
When you interact with a Cardano DEX today, the contract executes all calculations on-chain: it checks amounts, validates prices, and enforces invariants for \emph{every} order, consuming script budget. With ZK the flow changes completely:

\begin{enumerate}
  \item You build the transaction and perform all calculations \emph{off-chain} --- the amounts to receive, slippage bounds, fee splits.
  \item You attach to your transaction a ZK proof of those computations.
  \item The DEX contract checks only that the proof is valid. It does not re-execute the math.
\end{enumerate}

Because verification is orders of magnitude cheaper than re-computation, a single transaction can batch hundreds of orders. You see? ZK is being used to \textbf{compress computation}, not to hide anything.

\section{Core Concepts}

Before going further it is worth fixing the three properties that every honest ZK system must satisfy.

\begin{itemize}
  \item \textbf{Completeness.} If the statement is true and the prover is honest, the verifier will always be convinced.
  \item \textbf{Soundness.} If the statement is false, no cheating prover can convince the verifier (except with negligible probability).
  \item \textbf{Zero-Knowledge.} The verifier learns nothing beyond the truth of the statement. The proof reveals no information about the secret witness.
\end{itemize}

In every ZK protocol there are two parties:

\begin{itemize}
  \item \textbf{Prover} --- knows a secret (the \emph{witness}) and wants to convince the verifier without revealing it.
  \item \textbf{Verifier} --- checks the proof and decides whether to accept it.
\end{itemize}

\subsection*{Interactive vs.\ Non-Interactive Proofs}

Early ZK proofs were \emph{interactive}: the verifier sends random challenges and the prover responds, repeating rounds until the verifier is statistically convinced. This is elegant in theory but impractical on a blockchain where there is no live back-and-forth.

The \textbf{Fiat--Shamir heuristic} (1986) transforms interactive protocols into \emph{non-interactive} ones: the prover generates the verifier's challenges themselves by hashing the transcript so far. The resulting proof can be verified by anyone, at any time, without interaction --- exactly what on-chain verification requires.

\section{Types of Zero-Knowledge Proof Systems}

The ZK landscape has exploded in the past decade. Here is a practical map of the main families.

\begin{itemize}
  \item \textbf{zk-SNARKs} (Succinct Non-interactive Arguments of Knowledge) --- very small proofs and fast verification, but require a trusted setup ceremony. Groth16 and PLONK are the dominant constructions. Used in Zcash, Ethereum's zkSync and Scroll.
  \item \textbf{zk-STARKs} (Scalable Transparent Arguments of Knowledge) --- no trusted setup, post-quantum secure, but larger proof sizes. Used by StarkWare and StarkNet.
  \item \textbf{Bulletproofs} --- no trusted setup, compact range proofs, but verification is slower and scales linearly with circuit size. Used in Monero and Mimblewimble.
  \item \textbf{PLONK and universal SNARKs} --- use a single trusted setup that works for \emph{any} circuit up to a given size, removing the need for per-circuit ceremonies.
\end{itemize}

\begin{center}
\begin{tabular}{lccc}
  \hline
  \textbf{System}     & \textbf{Trusted Setup} & \textbf{Proof Size} & \textbf{Verification} \\
  \hline
  Groth16             & Per-circuit            & $\sim$200 bytes     & Very fast             \\
  PLONK               & Universal              & $\sim$500 bytes     & Fast                  \\
  STARKs              & None                   & $\sim$45--200 KB    & Moderate              \\
  Bulletproofs        & None                   & $\sim$1 KB          & Slow (linear)         \\
  \hline
\end{tabular}
\end{center}

\section{ZK on Cardano}

Cardano's \textbf{eUTxO model} is well-suited to ZK verification: a script validator receives the proof as part of the redeemer or datum and checks it in a single deterministic execution. Because Cardano contracts run off-chain before submission, the prover can generate the proof locally and attach it --- the on-chain script only verifies.

With \textbf{Plutus V3}, Cardano gained native cryptographic primitives for the BLS12-381 elliptic curve pairing:

\begin{itemize}
  \item \texttt{bls12\_381\_G1\_add}, \texttt{bls12\_381\_G2\_add} --- elliptic curve point addition in $\mathbb{G}_1$ and $\mathbb{G}_2$.
  \item \texttt{bls12\_381\_millerLoop}, \texttt{bls12\_381\_finalVerify} --- pairing operations needed to verify Groth16 proofs on-chain.
\end{itemize}

These primitives are also exposed in \textbf{Aiken}'s standard library under \texttt{aiken/crypto/bls12\_381}, making it practical to write on-chain Groth16 verifiers in a high-level language.

The goal of this book is to start from the very basics --- proving to someone that you know a value and can therefore spend a UTxO --- and build up to more complex applications: compressing off-chain computation, and eventually constructions similar to privacy protocols (studied here for academic purposes only).
